\section{Conclusion}
This paper introduces \workname, a Data efficient Contrastive Language-Image Pre-training paradigm. Our goal is to learn visual representations through the use of broader and scalable supervision. Specifically, instead of using the single image-text contrastive supervision, we fully exploit data potential through the use of (1) self-supervision within each modality; (2) multi-view supervision across modalities; (3) nearest-neighbor supervision from other similar pairs. Experimentally, \workname shows superior effectiveness and efficiency with different types of neural nets(CNN and ViT) and different amounts of data. We hope our work could bring insights about exploiting the multi-modal data.

% This paper introduces \workname, a new vision language pre-training paradigm. In this paradigm, utilizing more supervision from three parts: self-supervision within each modality, multi-view supervision across modalities, nearest-neighbor supervision from other similar pairs, our \workname only uses 88M image-text pairs to exceed CLIP with 400M pairs in zero-shot ability and pre-training ability. This shows that our approach was successful in the effectiveness and efficiency. While we provide the best ResNet50 and ViT-B/32 models, we also provide a small but very effective benchmark for academia and industry, which allows to try explorations related to vision language pre-training with less computing resources. We believe that due to our contributions, the field of vision language pre-training will develop rapidly.